\documentclass[10pt,a4paper]{amsart}
\usepackage[margin=19mm]{geometry}
\usepackage{amsmath,amssymb,mathtools}
\usepackage[scr=rsfs,cal=euler]{mathalfa}
\usepackage{xfrac}
\usepackage[unicode,hidelinks,bookmarks=false]{hyperref}
\newcommand{\ie}{i.e.\ }
\newcommand{\cf}{cf.\ }
\newcommand{\resp}{respectively\ }
\newcommand{\Subsection}{\subsection}
\providecommand{\nobreakdash}{\nobreak\hbox{-}}
\setlength{\emergencystretch}{2em}
\multlinegap0pt
\usepackage{bm}
\usepackage{bbm}
\usepackage{dsfont}
\usepackage{xspace}
\usepackage{cancel}
\usepackage{mathtools}
\usepackage[shortlabels]{enumitem}
\usepackage{amsthm}
\makeatletter
\@ifundefined{theorem}{\newtheorem{theorem}{Theorem}}{}
\@ifundefined{assumption}{\newtheorem{assumption}[theorem]{Assumption}}{}
\@ifundefined{proposition}{\newtheorem{proposition}[theorem]{Proposition}}{}
\makeatother
\newcommand{\Fcal}{\mathcal{F}}
\newcommand{\R}{\mathbb{R}}
\newcommand{\abs}[1]{\left\lvert#1\right\rvert}
\newcommand{\norm}[1]{\left\lVert#1\right\rVert}
\title[Neural Expectation Operators]{Neural Expectation Operators}
\author{Qian Qi}
\date{Source: arXiv:2507.10607v1 (CC BY 4.0)}
\begin{document}
\maketitle

\noindent\emph{Source and licence.} Neural Expectation Operators. Version arXiv:2507.10607v1 (CC BY 4.0), distributed under the Creative Commons Attribution 4.0 International licence, as recorded in the arXiv licence metadata for this exact version and shown on the abstract page https://arxiv.org/abs/2507.10607v1. Full rights notice: licenses/arxiv-2507.10607-CC-BY-4.0.txt.\par\smallskip

\noindent\textit{Editorial scope.} Counted unit: the Proposition whose source label is \texttt{prop:architectures} in the article (source0.tex lines 255--265, source label \texttt{prop:architectures}), delivered together with the article's own complete original proof of it (source0.tex lines 267--286, one \texttt{proof} environment of 20 lines). No mathematics is written, repaired or re-derived: statement and proof are the article's own text line for line. Required context retained uncounted: ass:f\_theta\_regularity\_general (lines 232--244). Every definition, display and label of those lines is retained unchanged. Omitted from this package and not redistributed: 1--231, 245--254, 266--266, 287--1047. Nothing inside a retained excerpt is omitted or altered: every retained body is delivered exactly as the article's lines are written, so no omission receipt of any kind accompanies this package. Citations: the retained excerpts carry no citation command at all, so no bibliography entry of the article is transcribed and no reference is redistributed. Reference adaptation: none was needed and none was made; every reference inside the delivered unit resolves inside it, so no unresolved reference or citation is produced. Printed slips kept literal: no printed word, symbol, hypothesis or number of the counted statement or of its proof was altered. Layout and class: the article's own class is \texttt{imsart} with packages including amssymb, amsfonts, lmodern, microtype, booktabs, enumitem, bm, xcolor, natbib, mathtools, graphicx, caption, subcaption, tabularx; the wrapper is the lane's pinned \texttt{amsart} editorial wrapper at 10pt on A4, which restates the theorem-like environments the retained text needs and adds the article's own macro definitions that the retained text needs (\texttt{\textbackslash Fcal}, \texttt{\textbackslash R}, \texttt{\textbackslash abs}, \texttt{\textbackslash norm}), with extra wrapper packages (bm, bbm, dsfont, xspace, cancel, mathtools, enumitem). The delivered numbering is the unit's own. Assets: no retained span contains a figure environment, a graphics-inclusion command, a PDF-page inclusion, a table or a TikZ picture; no image file is copied into this package, the package declares no asset dependency and no image file is redistributed with it. Licence: version arXiv:2507.10607v1 of this article is distributed under the Creative Commons Attribution 4.0 International licence, as recorded in the arXiv licence metadata for this exact version and shown on the abstract page https://arxiv.org/abs/2507.10607v1. A full rights notice is delivered at licenses/arxiv-2507.10607-CC-BY-4.0.txt. Characters: every retained excerpt is pure ASCII, so the delivered engine, XeLaTeX with its default Latin Modern text font, reports no missing character.
\begin{assumption}[Regularity and Growth of the NN Driver $f_\theta$]
\label{ass:f_theta_regularity_general}
The neural network function $f_\theta: [0,T] \times \Omega \times \R \times \R^{1 \times d} \to \R$ is $(\Fcal_t)$-progressively measurable via its dependence on $X_t$. For each fixed $\theta$, it satisfies:
\begin{enumerate}[label=(\roman*)]
    \item \textbf{Continuity:} The function $(t,x,y,z) \mapsto f_\theta(t,x,y,z)$ is continuous for each fixed $\omega \in \Omega$.
    \item \textbf{Quadratic Growth in $z$:} There exist constants $K > 0$, $\alpha > 0$ and $p \ge 1$ such that for all $(t,x,y,z)$:
    \[ \abs{f_\theta(t,x,y,z)} \le K(1 + \norm{x}^p + \abs{y}) + \frac{\alpha}{2}\norm{z}^2. \]
    \item \textbf{Uniform Local Lipschitz Condition in $y$:} For any $R>0$, there exists a constant $L_R > 0$ such that for all $(t,x,z)$ and all $y_1, y_2 \in [-R, R]$, we have
    \[ \abs{f_\theta(t,x,y_1,z) - f_\theta(t,x,y_2,z)} \le L_R\abs{y_1-y_2}. \]
    The uniformity is critical: the constant $L_R$ depends only on the radius $R$, not on $(t,x,z)$.
    \item \textbf{Monotonicity in $y$:} The function $y \mapsto f_\theta(t,x,y,z)$ is non-increasing for all fixed $(t,x,z)$.
\end{enumerate}
\end{assumption}

\begin{proposition}[Architectures for Key Structural Properties]
\label{prop:architectures}
The conditions of Assumption \ref{ass:f_theta_regularity_general} can be satisfied by the following designs.
\begin{enumerate}[label=(\alph*)]
    \item \textbf{Separable Architecture:} Let $f_\theta(t,x,y,z) = N_1(t,x,z) + N_2(y)$, where $N_1$ is any network satisfying the quadratic growth condition in $z$ and $N_2$ is a 1D network that is locally Lipschitz. This satisfies (iii). If $N_2$ is further constrained to be non-increasing, it also satisfies (iv).

    \item \textbf{Bounded Interaction Architecture:} Let $f_\theta(t,x,y,z) = N_1(t,x,z) + N_2(t,x,z) \cdot N_3(y)$, where $N_1$ satisfies the required growth, $N_3$ is a locally Lipschitz network, and the output of $N_2$ is bounded, i.e., $\sup_{t,x,z} \abs{N_2(t,x,z)} \le M_2 < \infty$. This architecture satisfies (iii).

    \item \textbf{Monotonicity by Construction:} Let $f_\theta$ be an $L$-layer feed-forward network with non-decreasing, continuously differentiable activation functions $\sigma_k$. Let the output layer be linear. The network output can be made non-increasing with respect to its input $y$ by enforcing the following sign constraints on its weights: (1) weights connecting $y$ to the first hidden layer are non-positive; (2) all weights in subsequent hidden layers ($k>1$) are non-negative; (3) weights in the final linear output layer are non-negative.
\end{enumerate}
\end{proposition}

\begin{proof}
\textbf{(a) Separable Architecture:} Let $y_1, y_2 \in [-R,R]$. Since $N_2$ is a 1D network with locally Lipschitz activations (e.g., ReLU, SiLU), it is locally Lipschitz. Thus, there exists a constant $L_{R,2}$ such that $\abs{N_2(y_1)-N_2(y_2)} \le L_{R,2}\abs{y_1-y_2}$. We verify the uniform local Lipschitz condition for $f_\theta$:
\begin{align*}
    \abs{f_\theta(t,x,y_1,z) - f_\theta(t,x,y_2,z)} &= \abs{(N_1(t,x,z) + N_2(y_1)) - (N_1(t,x,z) + N_2(y_2))} \\
    &= \abs{N_2(y_1) - N_2(y_2)} \le L_{R,2}\abs{y_1-y_2}.
\end{align*}
The constant $L_{R,2}$ depends only on $R$, so (iii) is satisfied. Monotonicity (iv) follows directly if $N_2$ is chosen to be non-increasing.

\textbf{(b) Bounded Interaction Architecture:} Let $y_1, y_2 \in [-R,R]$. $N_3$ is locally Lipschitz, so there exists $L_{R,3}$ such that $\abs{N_3(y_1)-N_3(y_2)} \le L_{R,3}\abs{y_1-y_2}$. The network $N_2$ is bounded by $M_2$.
\begin{align*}
    \abs{f_\theta(t,x,y_1,z) - f_\theta(t,x,y_2,z)} &= \abs{N_2(t,x,z)(N_3(y_1) - N_3(y_2))} \\
    &\le \abs{N_2(t,x,z)} \cdot \abs{N_3(y_1) - N_3(y_2)} \le M_2 \cdot L_{R,3} \abs{y_1 - y_2}.
\end{align*}
The constant $L_R \coloneqq M_2 L_{R,3}$ depends only on $R$, so (iii) holds.

\textbf{(c) Monotonicity by Construction:} Let the network input be the vector $u \coloneqq (t,x,y,z)$. Let $a^{(0)} \coloneqq u$. For layer $k \in \{1, \dots, L-1\}$, the pre-activation is $z^{(k)} = W^{(k)} a^{(k-1)} + b^{(k)}$ and the activation is $a^{(k)} = \sigma_k(z^{(k)})$. The final output is $f_\theta = W^{(L)}a^{(L-1)} + b^{(L)}$. We show by induction that $\frac{\partial a_i^{(k)}}{\partial y} \le 0$ for all neurons $i$ and layers $k \in \{1, \dots, L-1\}$.
\textit{Base Case (k=1):} The derivative of the $i$-th neuron's activation in the first hidden layer with respect to $y$ is $\frac{\partial a_i^{(1)}}{\partial y} = \sigma_1'(z_i^{(1)}) \frac{\partial z_i^{(1)}}{\partial y} = \sigma_1'(z_i^{(1)}) W_{iy}^{(1)}$. Since $\sigma_1' \ge 0$ and we constrain $W_{iy}^{(1)} \le 0$, the derivative is non-positive.
\textit{Inductive Step:} Assume for layer $k-1$ that $\frac{\partial a_j^{(k-1)}}{\partial y} \le 0$ for all neurons $j$. For a neuron $i$ in layer $k$, the chain rule gives $\frac{\partial a_i^{(k)}}{\partial y} = \sigma_k'(z_i^{(k)}) \sum_j W_{ij}^{(k)} \frac{\partial a_j^{(k-1)}}{\partial y}$. Since $\sigma_k' \ge 0$, we constrain $W_{ij}^{(k)} \ge 0$ for $k>1$, and by the induction hypothesis, $\frac{\partial a_j^{(k-1)}}{\partial y} \le 0$, each term in the sum is non-positive. Thus, the sum is non-positive, and so is $\frac{\partial a_i^{(k)}}{\partial y}$.
\textit{Conclusion:} The derivative of the final network output with respect to $y$ is $\frac{\partial f_\theta}{\partial y} = \sum_j W_{j}^{(L)} \frac{\partial a_j^{(L-1)}}{\partial y}$. Since the final layer weights $W_j^{(L)}$ are constrained to be non-negative and we have shown $\frac{\partial a_j^{(L-1)}}{\partial y} \le 0$, the final derivative is non-positive. This proves that $f_\theta$ is non-increasing with respect to $y$.
\end{proof}

\end{document}
